%%% BEGIN fig:pdf-operator
\begin{figure}[htbp]
\centering
\tikzset{t1 hadron/.style={draw,solid,line width=1.5pt,line cap=butt}}
\begin{tikzpicture}[fd]
  \coordinate (ol) at (-1.3,2.3); \coordinate (or) at (1.3,2.3);
  \coordinate (bl) at (-1.5,0);   \coordinate (br) at (1.5,0);
  % remnant X crossing the cut (drawn first; the blobs cover the ends)
  \foreach \y in {0.24,0,-0.24} \draw[plain] (bl) ++(0,\y) -- ++(3,0);
  \draw[quark] (-1.3,0.43) -- (ol);
  \draw[wilson=0.27] (ol) -- (or);
  \draw[quark] (or) -- (1.3,0.43);
  \draw[t1 hadron] (-3.25,0) -- (bl);
  \draw[t1 hadron] (br) -- (3.25,0);
  \draw[hadron blob] (bl) ellipse[x radius=0.82,y radius=0.45];
  \draw[hadron blob] (br) ellipse[x radius=0.82,y radius=0.45];
  \node[opvertex] at (ol) {}; \node[opvertex] at (or) {};
  \draw[cut] (0,-0.85) -- (0,2.9);
  % labels
  \path[mom={$k$}] (-1.3,0.45) -- (ol);
  \path[mom={$k$}] (or) -- (1.3,0.45);
  \node[left,font=\small] at (-1.6,1.0) {$k^+=\xi P^+$};
  \node[above left,inner sep=3pt] at (ol) {$\psi(0)$};
  \node[above right,inner sep=3pt] at (or) {$\bar\psi(w)$};
  \node[above,font=\small,inner sep=4pt] at (-0.65,2.3) {$W$};
  \path[mom={$P$},mom sep=0.3cm] (-3.25,0) -- (-2.32,0);
  \path[mom={$P$},mom sep=0.3cm] (2.32,0) -- (3.25,0);
  \node[font=\small] at (0.42,-0.62) {$X$};
\end{tikzpicture}
\caption{夸克分布函数的割图表示。插入完备态并完成 $w^-$ 积分，
$f_{q/H}(\xi)=\sum_X\tfrac12\,\delta\big(\xi P^+-(P-p_X)^+\big)\,\langle P|\bar\psi(0)W[0,\infty n]|X\rangle\,\gamma^+\,\langle X|W[\infty n,0]\psi(0)|P\rangle$，
其中 $n$ 是满足 $n\cdot v=v^+$ 的光锥矢量（$w=w^-n$）。割线左侧为振幅、右侧为共轭振幅，$X$ 为穿过割线的剩余部分。
两个 $\otimes$ 为非定域算符 $\bar\psi(w)\gamma^+W[w,0]\psi(0)$ 的两端，算符顶点的 Feynman 规则为 $\tfrac12\gamma^+\delta(k^+-\xi P^+)$
（对 $\int\dd^4k/(2\pi)^4$；Wilson 线上没有胶子附着时）；连接两端的双线为 Wilson 线 $W[w,0]=W[w,\infty n]\,W[\infty n,0]$，
被末态割线一分为二。箭头表示费米子数（色）流向。}
\label{fig:pdf-operator}
\end{figure}
%%% END fig:pdf-operator
